\section{Related Work and Positioning}
\label{sec:related}
Time-distributed FFT execution has a substantial history. Lo and Lee construct
butterfly work during acquisition and also combine scheduling with recursive
reuse across shifted frames~\cite{lo1998,lo1999,lomoving1999}. In convolution,
Bor\ss{} distributes partition work and staggers filter schedules to control
their combined load~\cite{borss2009}, while Hurchalla distributes transform
work across incoming blocks on one thread~\cite{hurchalla2010}. Battenberg and
Avi\v{z}ienis schedule forward transforms, spectral products, and inverse
transforms across callbacks, use FFTW for the leaf transforms, and compare
cooperative with preemptive execution~\cite{battenberg2011}. Complete
processing chains and optimized transform kernels thus both have precedents
within cooperative scheduling. Wefers treats manual preemption and transform
subdivision more broadly~\cite{wefers2015}, making clear that suspension
granularity is an implementation choice rather than an inherent advantage of
butterfly-level execution. Liu, Yan, and Zou study time-distributed FFT and
inverse FFT execution for online control~\cite{liu2017}. We therefore claim
neither transform suspension nor cooperative processing chains as new
principles.

For spectral visualization, Pr\r{u}\v{s}a and Holighaus use a worker thread
and distinguish computation, polling, and repaint delay~\cite{prusa2016}.
Their filter-bank representation differs from the fixed-window FFT considered
here, but their separation of execution context from visible-output latency
is directly relevant. Sliding, hopping, and feedforward transforms offer a
different dimension of choice: reuse of work across overlapping
windows~\cite{lyons2021,park2014,garrido2016}, which frequency-domain kernels
can extend to tapered windows~\cite{rafii2018,eleftheriadis2023}. Overlap
reuse and time distribution are complementary rather than mutually exclusive.
Our implementation selects a complete window and starts a fresh transform for
each frame.

Partitioned convolution makes the tradeoff between latency and cost explicit
in a different application: Gardner combines fast convolution with low
input/output delay, Garcia optimizes nonuniform filter partitions for a
specified delay, and Wefers and Vorl\"ander optimize uniform FFT partition
sizes~\cite{gardner1995,garcia2002,wefers2011}. These precedents motivate
treating execution granularity as a design variable; this paper introduces no
new convolution algorithm.

\begin{table}[htbp]
\centering\small
\begin{tabular}{@{}p{.26\linewidth}p{.31\linewidth}p{.35\linewidth}@{}}
\toprule
Prior approach & Established contribution & Focus of this study \\
\midrule
Time-distributed FFTs~\cite{hurchalla2010,lo1999} & Transform work spread over acquisition or processing intervals & Complete analysis tail, exact-hop producer cadence, and retained selected input. \\
Cooperative convolution~\cite{battenberg2011} & Scheduled chains with optimized native leaves & Analysis-specific magnitude/smoothing/display output and matched placement controls. \\
Transform subdivision~\cite{wefers2015} & Granularity as an implementation choice & Measured batch, whole-native hybrid, and butterfly paths; native leaves remain unmeasured. \\
Worker visualization~\cite{prusa2016} & Separate computation, polling, and repaint & Engine-thread publication ownership and bounded logical consumer-age observations. \\
\bottomrule
\end{tabular}
\caption{Contribution relative to closely related execution models. The
right column states this study's emphasis, not an exhaustive absence claim
about prior implementations.}
\label{tab:contribution-delta}
\end{table}


The present study specializes that scheduling perspective to windowed
analysis (Table~\ref{tab:contribution-delta}): preparation and the
magnitude, smoothing, and display stages belong inside the contract;
exact-hop publication is independent of rounding in the work quota; and
matched controls distinguish placement from kernel choice. FFTW, Rack's
PFFFT wrapper, and Apple's Accelerate/vDSP provide practical native
alternatives~\cite{frigo2005,pffft,vdsp}. A headless evaluation inside Rack's
engine then tests whether differences in local concentration survive the host
boundary. Appendix~\ref{app:related-work} gives a broader comparison with the
literature.
